\chapter{Purpose and Scope}

\section{Purpose}

This document defines the first NavKit IMU emulator contract. It is a
code-ready specification for turning truth angular-rate and specific-force
signals into timestamped IMU increments consumed by
\texttt{navigator\_ecef\_v1}. It extracts the useful triad error-model content
from the broader navigation reference while using cleaner implementation-facing
notation \cite{groves2013,titterton2004,grewal2020}.

\begin{scopebox}
The goal is not to prove every equation from first principles. The goal is to
state every equation and convention needed to implement, configure, and test
the IMU emulator without inventing math in C++.
\end{scopebox}

\section{V1 Scope}

Version 1 assumes:
\begin{itemize}
    \item one IMU;
    \item IMU frame collocated with the body frame and navigation state point;
    \item truth angular rate and specific force already resolved in body;
    \item no IMU lever arm;
    \item no multiple-IMU voting or fusion;
    \item no online estimation of scale-factor, misalignment, or
    non-orthogonality states;
    \item deterministic seeded random draws for repeatable simulation.
\end{itemize}

The emulator may still model fixed calibration errors, time-varying biases,
white measurement noise, and quantization because those are needed to exercise
the downstream Navigator realistically.

\begin{implbox}
Trajectory providers supply ECEF truth states only. They do not supply IMU truth
directly. The IMU emulator derives ideal inertial IMU increments from successive
ECEF truth states, including the Earth-rotation correction needed to convert
body-with-respect-to-ECEF attitude changes into body-with-respect-to-inertial
gyro increments.
\end{implbox}

\section{Deferred Generalizations}

The following are explicitly deferred:
\begin{itemize}
    \item IMU lever-arm correction;
    \item multiple IMUs;
    \item temperature-dependent error models;
    \item vibration/colored-noise models beyond simple bias random walk;
    \item online calibration-state estimation;
    \item sensor-frame mounting distinct from the body frame in the
    mechanization itself.
\end{itemize}

Those features can be added after the single body-collocated IMU path is tested.
